%% Teste mínimo para verificação básica de compilação
%% Este arquivo testa apenas os elementos essenciais da classe UnB-PPGT

\documentclass[mestrado,12pt]{UnB-PPGT}

%% Metadados mínimos obrigatórios
\titulo{Teste Mínimo de Compilação}
\autor{Sistema}{de Testes}
\orientador{Prof. Dr. Orientador Teste}{Universidade de Brasília}
\diamesano{1}{1}{2024}
\coordenador{Prof. Dr. Coordenador Teste}{Universidade de Brasília}

%% Banca mínima (3 membros)
\membrobancafunc{Prof. Dr. Orientador Teste}{Universidade de Brasília}{Orientador}
\membrobancafunc{Prof. Dr. Membro 1}{Universidade de Brasília}{Examinador}
\membrobancafunc{Prof. Dr. Membro 2}{Universidade Federal}{Examinador}

%% Palavras-chave obrigatórias
\palavraschave{teste, compilação, latex}{teste, compilação, latex}
\keywords{test, compilation, latex}{test, compilation, latex}

\begin{document}

\pretextual

%% Páginas pré-textuais básicas
\capa
\folhaderosto
\folhadeaprovacao

%% Resumo e abstract mínimos
\begin{resumo}
Este é um teste mínimo de compilação da classe UnB-PPGT. O objetivo é verificar se todos os elementos essenciais da classe funcionam corretamente com configurações básicas.
\end{resumo}

\begin{abstract}
This is a minimal compilation test for the UnB-PPGT class. The objective is to verify that all essential elements of the class work correctly with basic configurations.
\end{abstract}

\textual

%% Conteúdo textual mínimo
\chapter{Introdução}
Este é um capítulo de teste com conteúdo mínimo para verificar a formatação básica.

\section{Seção de Teste}
Esta seção contém apenas texto simples para validar a estrutura do documento.

\chapter{Conclusão}
Esta é a conclusão do teste mínimo.

\postextual

\end{document}